\section{Results}
\label{sec:results}

\subsection{Main Results}

Table~\ref{tab:main_results} presents the mean \asr{}, \bsr{}, and \fpr{} across all 4 models $\times$ 3 runs for each defense strategy. All 11 defenses are compared on the same 250-case core benchmark (155 attack + 95 benign). For D8--D10, which were originally evaluated on the full 565-case benchmark, we extract the matched 250-case subset to ensure a fair cross-defense comparison. Statistical significance is assessed using McNemar's test with Bonferroni correction. Full 565-case results for D8--D10 appear in Appendix~\ref{sec:extended_results}.

\begin{table*}[t]
\centering
\caption{Unified comparison of all 11 defenses on the 250-case core benchmark (155 attack + 95 benign), averaged across 4 frontier models $\times$ 3 runs. D8--D10 results computed on the matched 250-case subset of their 565-case evaluations. 95\% Wilson CIs in brackets. Best non-baseline values in \textbf{bold}. $\Delta$ASR = relative change from D0 baseline.}
\label{tab:main_results}
\small
\begin{tabular}{ll rrrr}
\toprule
\textbf{ID} & \textbf{Defense} & \textbf{ASR$\downarrow$} & \textbf{BSR$\uparrow$} & \textbf{FPR$\downarrow$} & \textbf{$\Delta$ASR} \\
\midrule
\multicolumn{6}{l}{\textit{No defense}} \\
D0 & Baseline          & .413\tiny$\pm$.02 & .912\tiny$\pm$.02 & .088 & ---      \\
\midrule
\multicolumn{6}{l}{\textit{Prompt-level defenses}} \\
D5 & Sandwich          & \textbf{.010}\tiny$\pm$.00 & \textbf{.934}\tiny$\pm$.02 & \textbf{.066} & $-$97.5\%$^{***}$ \\
D1 & Spotlighting      & .020\tiny$\pm$.01 & .913\tiny$\pm$.02 & .087 & $-$95.1\%$^{***}$ \\
D7 & Input Classifier  & .430\tiny$\pm$.02 & .922\tiny$\pm$.02 & .078 & +4.0\%   \\
\midrule
\multicolumn{6}{l}{\textit{Tool-gating defenses}} \\
D8 & Semantic Firewall & .107\tiny$\pm$.01 & .383\tiny$\pm$.02 & .617 & $-$74.0\%$^{***}$ \\
D2 & Policy Gate       & .307\tiny$\pm$.02 & .762\tiny$\pm$.02 & .238 & $-$25.7\%$^{**}$  \\
D4 & Re-execution      & .322\tiny$\pm$.02 & .779\tiny$\pm$.02 & .221 & $-$22.2\%$^{*}$  \\
D3 & Task Alignment    & .406\tiny$\pm$.02 & .916\tiny$\pm$.02 & .084 & $-$1.8\%  \\
D9 & Dual-LLM          & .263\tiny$\pm$.04 & .578\tiny$\pm$.04 & .422 & $-$36.5\%$^{*}$ \\
\midrule
\multicolumn{6}{l}{\textit{Content-level defense}} \\
D6 & Output Filter     & .379\tiny$\pm$.02 & .902\tiny$\pm$.02 & .098 & $-$8.3\%  \\
\midrule
\multicolumn{6}{l}{\textit{Multi-signal defense (ours)}} \\
D10 & CIV              & .089\tiny$\pm$.01 & .821\tiny$\pm$.02 & .179 & $-$78.4\%$^{***}$ \\
\bottomrule
\multicolumn{6}{l}{\footnotesize $^{***}p<0.001$, $^{**}p<0.01$, $^{*}p<0.05$ McNemar's test vs.\ D0 (Bonferroni-corrected across models).}
\end{tabular}
\end{table*}


\paragraph{Finding 1: Prompt-Level Defenses Dominate.}
\defense{D5} (Sandwich) achieves the lowest \asr{} (1.0\%) with the highest \bsr{} (93.4\%), confirming that goal re-anchoring is highly effective.
\defense{D1} (Spotlighting) ranks second with 2.0\% \asr{}.
Both defenses reduce \asr{} by $>$95\% relative to the D0 baseline while preserving $>$91\% \bsr{}.
This demonstrates that prompt-layer intervention achieves the most favorable security--utility trade-off, challenging the intuition that more complex mechanisms are needed.

\paragraph{Finding 2: Tool-Gating Defenses Trade Utility for Security.}
Tool-gating defenses split into two clusters.
\emph{Aggressive gates}---\defense{D8} (Semantic Firewall, 74\% ASR reduction but only 38.3\% BSR) and \defense{D9} (Dual-LLM, 36.5\% reduction, 57.8\% BSR)---achieve strong attack suppression but suffer from catastrophic over-blocking (\fpr{} $>$42\%), rendering them impractical as standalone defenses.
\emph{Moderate gates}---\defense{D2} (Policy Gate, 25.7\% reduction) and \defense{D4} (Re-execution, 22.2\%)---show more balanced trade-offs with \fpr{} $\sim$22--24\%.
\defense{D3} (Task Alignment) provides negligible improvement ($-$1.8\%), suggesting that LLM-based consistency checks alone are insufficient to detect injections.

\paragraph{Finding 3: CIV Provides a Complementary Mechanism.}
\defense{D10} (\civ{}) achieves 8.9\% \asr{}---a 78.4\% reduction from baseline---while maintaining 82.1\% \bsr{} on the 250-case core benchmark.
On the full 565-case benchmark, \bsr{} drops to 50.3\%; however, the D0 baseline also drops to 53.1\%, indicating that this decline is largely attributable to mock-environment limitations on complex multi-tool workflows rather than CIV-specific false positives (see \S\ref{sec:limitations} and Appendix~\ref{sec:extended_results}).
Nevertheless, operating through tool-call gating rather than prompt modification, CIV demonstrates that information-flow-based approaches provide a complementary defense mechanism when prompt-level defenses are already deployed.

\paragraph{Finding 4: Model Robustness Varies Widely.}
Claude~4.5 Sonnet exhibits the lowest baseline \asr{} (11.6\%), approximately 4--5$\times$ more robust than GPT-4o (49.2\%), DeepSeek~V3.1 (50.8\%), and Gemini~2.5 Flash (53.8\%).
This large gap suggests that instruction-following alignment varies substantially across model families.
However, all models benefit from effective defenses: D5 achieves $<$1.3\% \asr{} across all four models (Figure~\ref{fig:model_comparison}, Appendix).

\subsection{CIV Ablation Study}
\label{sec:ablation}

Table~\ref{tab:civ_ablation} reports ablation results for \civ{}'s three signal components (GPT-4o, 250-case core benchmark).

\begin{table}[t]
\centering
\small
\caption{CIV ablation study on the 250-case core benchmark (GPT-4o, single run). Top: CIV~v1 component ablation. Bottom: CIV~v2 (main table) incorporating ablation insights. $\Delta$ASR relative to v1-full.}
\label{tab:civ_ablation}
\resizebox{\columnwidth}{!}{%
\begin{tabular}{l rr rr}
\toprule
Config & ASR$\downarrow$ & BSR$\uparrow$ & $\Delta$ASR & $\Delta$BSR \\
\midrule
\textbf{full}                 & 0.123 & 0.853 & ---        & ---      \\
\midrule
\multicolumn{5}{l}{\textit{Leave-one-out}} \\
no\_provenance                & 0.265 & 0.905 & +115\%     & +6\%     \\
no\_fingerprint               & 0.116 & 0.747 & $-$6\%     & $-$12\%  \\
no\_counterfactual            & 0.058 & 0.821 & $-$53\%    & $-$4\%   \\
\midrule
\multicolumn{5}{l}{\textit{Single-signal}} \\
provenance\_only              & 0.097 & 0.832 & $-$21\%    & $-$2\%   \\
fingerprint\_only             & 0.361 & 0.821 & +193\%     & $-$4\%   \\
counterfactual\_only          & 0.226 & 0.811 & +84\%      & $-$5\%   \\
\midrule
\multicolumn{5}{l}{\textit{CIV v2: incorporating ablation insights (main table)}} \\
\textbf{v2} (embed + plan-dev)   & \textbf{0.089} & 0.821 & $-$28\%    & $-$4\%   \\
\bottomrule
\end{tabular}%
}
\begin{flushleft}
\scriptsize
\textbf{Key findings}: (1)~Provenance is the most critical signal (+115\% ASR when removed); single-signal provenance (9.7\% ASR) nearly matches v1-full.
(2)~Counterfactual reasoning is net-harmful ($-$53\% ASR when removed), motivating its replacement in v2.
(3)~Fingerprint (tool affinity) is critical for utility ($-$12\% BSR when removed), preserved in v2 via embedding-based compatibility.
(4)~CIV~v2 replaces counterfactual with plan deviation detection and upgrades fingerprint to embedding similarity, achieving 8.9\% ASR ($-$28\% vs.\ v1) while maintaining comparable BSR.
\end{flushleft}
\end{table}


\paragraph{Provenance Is the Core Signal.}
Removing provenance causes the largest \asr{} increase (+115\%), confirming entity provenance tracking as the dominant attack detection mechanism.
Conversely, provenance alone achieves 9.7\% \asr{}---nearly matching the full system---suggesting it captures the primary attack signal.

\paragraph{Counterfactual Reasoning Is Counterproductive (v1 Finding).}
Surprisingly, removing the heuristic counterfactual signal from CIV~v1 \emph{improves} security ($-$53\% \asr{}).
The counterfactual heuristic is too permissive, inadvertently allowing attacks that provenance would block.
This finding directly motivated the v1$\to$v2 redesign: CIV~v2 replaces this signal with plan deviation detection, achieving 8.9\% ASR ($-$28\% vs.\ v1-full).

\paragraph{Tool Affinity Supports Utility.}
The fingerprint (tool affinity) signal is most important for maintaining utility: removing it causes the largest BSR drop ($-$12\%).
It helps CIV distinguish legitimate tool calls by modeling goal--tool compatibility.

\paragraph{Weight Robustness.}
To verify that \civ{}'s performance is not overfit to a specific weight configuration, we evaluate five weight settings: the default (0.45/0.30/0.25), equal weights (0.33/0.33/0.33), and three single-signal-heavy variants (provenance 0.60, compatibility 0.50, or plan 0.50).
All five configurations produce identical ASR, BSR, and FPR on our evaluation set, indicating that \civ{}'s effectiveness is driven by the signal design rather than weight tuning.
Full results appear in Appendix~\ref{app:weight_sensitivity}.

\subsection{Cross-Benchmark Generalization}
\label{sec:cross_benchmark}

We tested selected defenses on InjecAgent \citep{zhan2024injecagent} and AgentDojo \citep{debenedetti2024agentdojo} using adapter-based format conversion.
All defenses---including the D0 baseline---achieve 0\% \asr{} on both external benchmarks, revealing that \textbf{attack success judgments are tightly coupled to each benchmark's evaluation protocol} and cannot be transferred via format conversion alone.
The mismatch arises from differences in forbidden-action semantics, task-tool coupling, and environment dynamics.
This finding implies that \textbf{cross-benchmark defense comparisons require benchmark-native evaluation}, an open challenge for the community.
